\documentclass[10.5pt]{article}
\usepackage{amsmath, amsfonts, amssymb,amsthm}
\usepackage[includeheadfoot]{geometry} % For page dimensions
\usepackage{fancyhdr}
\usepackage{enumerate} % For custom lists
\usepackage{tikz-cd}
\usepackage{graphicx}

\fancyhf{}
\lhead{Engel's Theorem}
\rhead{Tighe McAsey}
\pagestyle{fancy}

% Page dimensions
\geometry{a4paper, margin=1in}

\theoremstyle{definition}
\newtheorem{pb}{Exercise}
\newtheorem{definition}{Definition}
\newtheorem{proposition}{Proposition}
\newtheorem{remark}{Remark}
\newtheorem{example}{Example}
\newtheorem{theorem}{Theorem}
\numberwithin{example}{subsection}
\usepackage{tikz-cd, stackengine}

% Commands:

\newcommand{\set}[1]{\left\{#1\right\}}
\newcommand{\gen}[1]{\langle#1\rangle}
\newcommand{\abs}[1]{\left\vert#1\right\vert}
\newcommand{\norm}[1]{\lvert\lvert#1\rvert\rvert}
\newcommand{\tand}{\text{ and }}
\newcommand{\tor}{\text{ or }}
\newcommand{\pd}{\frac{\partial}{\partial x_j}}
\newcommand{\px}{\frac{\partial}{\partial x}}
\newcommand{\py}{\frac{\partial}{\partial y}}
\newcommand{\pz}{\frac{\partial}{\partial z}}
\newcommand{\ppx}{\frac{\partial^2}{\partial x^2}}
\newcommand{\ppy}{\frac{\partial^2}{\partial y^2}}
\newcommand{\ppz}{\frac{\partial^2}{\partial z^2}}
\newcommand{\hess}{\operatorname{Hess}}
\newcommand{\inv}{\operatorname{inv}}
\newcommand{\spl}{\operatorname{\mathfrak{sl}}}
\newcommand{\spg}{\mathfrak{g}}
\newcommand{\sph}{\mathfrak{h}}
\newcommand{\spu}{\operatorname{\mathfrak{su}}}
\newcommand{\spo}{\operatorname{\mathfrak{so}}}

\begin{document}
    \begin{remark}
        It was given as an exercise in class to define the semisimple and nilpotent parts of an element \(x \in \mathfrak{g}\), we do that here.

        \begin{proof}
            The proof uses two results:
            \begin{proposition}[Humphreys Lemma 4.2B]\label{prop}
                Let \(\mathfrak{U}\) be a finite dimensional \(k\)-Algebra. Then for \(x \in \mathfrak{U}\) the jordan decompoisition \(x = x_s + x_n\) in \(\operatorname{End}(\mathfrak{U}) \supset \operatorname{Der}(\mathfrak{U})\) satisfies \(x_n, x_s \in \operatorname{Der}(\mathfrak{U})\).
            \end{proposition}

            \begin{proposition}[Humphreys Theorem 5.3]\label{prop2}
                If \(\mathfrak{g}\) is semisimple then \(\operatorname{ad}(\mathfrak{g}) = \operatorname{Der}(\mathfrak{g})\).
            \end{proposition}

            We are now ready to define the semisimple and nilpotent part of an element \(x\). Recall that \(\operatorname{ad}: \mathfrak{g} \to \mathfrak{gl}(\mathfrak{g})\), then for \(x \in \mathfrak{g}\) it follows that \(\operatorname{ad}(x)\) has a Jordan decomposition. Now since \(\mathfrak{g}\) is semisimple, we know that \(\mathfrak{z}(\mathfrak{g}) \subset \operatorname{Rad}(\mathfrak{g}) = 0\) and thus \(\operatorname{ad}: \mathfrak{g} \overset{\cong}{\longrightarrow} \operatorname{ad}(\mathfrak{gl}(\mathfrak{g}))\). Thus we define the Jordan decomposition as follows
            \begin{align}
                x_s = \operatorname{ad}^{-1}(\operatorname{ad}(x)_s) \tand x_n = \operatorname{ad}^{-1}(\operatorname{ad}(x)_n)
            \end{align}
            The fact that \(\operatorname{ad}(x)_s,\operatorname{ad}(x)_n \in \operatorname{ad}(\mathfrak{g})\) follows from Proposition \ref{prop} and Proposition \ref{prop2}.
        \end{proof}
    \end{remark}

    \begin{pb}[Humphrey's 5.8]
        Let \(\mathfrak{g}\) be a Lie algebra and \(\bigoplus I_j\) be its decomposition into simple ideals. Let \(\pi_j: \mathfrak{g} \to I_j\) be the projection, then show that
        \begin{align}
            x_n = \sum_j \pi_j(x)_n \tand x_s = \sum_j \pi_j(x)_s
        \end{align}

        Before beginning the proof it is worth noting that this decomposition uses the forward implication of \(\mathfrak{g}\) is semisimple \(\iff \kappa\) is nondegerate. This requires that the base field \(k\) have characteristic zero since the proof uses Cartan's criterion which relies on characteristic zero. A number of the results stated in defining the semisimple nilpotent decomposition also require certain niceties of the base field, for simplicity I assume that \(k\) is an algebraically closed field of characteristic zero, all of the stated propositions hold in this context.

        \begin{proof} Since \(\mathfrak{g}\) is semisimple we find that for \(i \neq j\) we have \([I_j,I_i] \subset I_j \cap I_i = 0\) so that
        \begin{align}
            [\operatorname{ad}(\pi_j(x)),\operatorname{ad}(\pi_i(x))] = \operatorname{ad}([\pi_j(x),\pi_i(x)]) = \operatorname{ad}(0) = 0
        \end{align}
        So all of the \(\operatorname{ad}(\pi_j(x))\) commute. We use a fact from linear algebra that the sum of commuting semisimple (resp. nilpotent) elements is semisimple (resp. nilpotent). It follows that
        \begin{align}
            \operatorname{ad}(x) = \sum_j \operatorname{ad}(\pi_j(x)) = \sum_j \operatorname{ad}(\pi_j(x))_s + \operatorname{ad}(\pi_j(x))_n = \underbrace{\sum_j \operatorname{ad}(\pi_j(x))_s}_{\text{semisimple}} + \underbrace{\sum_j \operatorname{ad}(\pi_j(x))_n}_{\text{nilpotent}}
        \end{align}
        Then by uniqueness of the Jordan decomposition, we find that
        \begin{align}
            \operatorname{ad}(x)_s = \sum_j \operatorname{ad}(\pi_j(x))_s \tand \operatorname{ad}(x)_n = \sum_j \operatorname{ad}(\pi_j(x))_n
        \end{align}
        It follows that
        \begin{align}
            x_s = \operatorname{ad}^{-1}(\operatorname{ad}(x)_s) = \operatorname{ad}^{-1}\left(\sum_j \operatorname{ad}(\pi_j(x))_s\right) = \sum_j \operatorname{ad}^{-1}(\operatorname{ad}(\pi_j(x))_s) = \sum_j \pi_j(x)_s \\
            x_n = \operatorname{ad}^{-1}(\operatorname{ad}(x)_n) = \operatorname{ad}^{-1}\left(\sum_j \operatorname{ad}(\pi_j(x))_n\right) = \sum_j \operatorname{ad}^{-1}(\operatorname{ad}(\pi_j(x))_n) = \sum_j \pi_j(x)_n
        \end{align}
    \end{proof}
    \end{pb}
\end{document}